% =====================================================================
%  ch06b —— 原书 p.80–95（PDF 89–104）
%  第 6 章 完全可积系统：6.4 R 矩阵与经典 Yang–Baxter 方程 / 6.5 AKS 定理
%                        6.6 因子化定理 / 6.7 Toda 格子再论
%  本章最后一块（\chapterreview 在本文件末尾）
% =====================================================================

% [原文 p.80]
\begin{equation}
  A = L + \dot{R}R^{-1}
  \tag{6.3.33}
\end{equation}
由此得到 (6.3.33)。这里，$R$ 为上三角矩阵，故 $\dot{R}$ 以及 $\dot{R}R^{-1}$ 也是上三角矩阵，因此 $A$ 与 $L$ 除对角分量以外的下三角部分一致。由这一点与 $A$ 的反称性可知 $A=M$，定理 6.3.1 证毕。\ifdefined\UEDAdraft\else\end{proof}\fi

\section{R 矩阵与经典 Yang--Baxter 方程}

对 Lie 代数 $(\mathfrak{g},[-,-])$ 而言，若 $\mathfrak{g}$ 的线性自同态 $R\in\End\mathfrak{g}$ 使得由
\begin{equation}
  [X,Y]_R = \frac{1}{2}\bigl([RX,Y] + [X,RY]\bigr)
  \tag{6.4.1}
\end{equation}
定义的双线性映射 $[-,-]_R:\mathfrak{g}\otimes\mathfrak{g}\to\mathfrak{g}$ 在 $\mathfrak{g}$ 上确定一个 Lie 代数结构，则称 $R$ 为 \textbf{R 矩阵}（R-matrix）。

对任意 $R\in\End\mathfrak{g}$，由 (6.4.1) 定义的 $[-,-]_R$ 显然具有交错性。给定 $R\in\End\mathfrak{g}$，定义线性映射 $B_R:\mathfrak{g}\otimes\mathfrak{g}\to\mathfrak{g}$ 为
\begin{equation}
  B_R(X,Y) = [RX,RY] - R\bigl([RX,Y] + [X,RY]\bigr)
  \tag{6.4.2}
\end{equation}

% [原文 p.80] （续）
\[
\begin{aligned}
4[[X,Y]_R,Z]_R
  &= 2\bigl[[RX,Y] + [X,RY],Z\bigr]_R \\
  &= 2[[RX,Y],Z]_R + 2[[X,RY],Z]_R \\
  &= [R[RX,Y],Z] + [[RX,Y],RZ] \\
  &\qquad + [R[X,RY],Z] + [[X,RY],RZ] \\
  &= [R[RX,Y],Z] + [R[X,RY],Z] \\
  &\qquad + [[RX,Y],RZ] + [[X,RY],RZ] \\
  &= [R[RX,Y],Z] + [R[X,RY],Z] \\
  &\qquad - [[Y,RZ],RX] - [[RZ,RX],Y] + [[X,RY],RZ] \\
  &= -[[RZ,RX],Y] + [R[RX,Y],Z] + [R[X,RY],Z] \\
  &\qquad - [[Y,RZ],RX] + [[X,RY],RZ]
\end{aligned}
\]
% [Tutor Note: 首行左端系数 4 来自 (6.4.1) 中的 1/2；右端于是化为 2[[RX,Y]+[X,RY],Z]_R，逐项展开即得下列各式。]

对 $(X,Y,Z)$ 作循环置换，将所得三式相加，则 $[-,-]_R$ 的 Jacobi 恒等式被改写为
\begin{equation}
  [B_R(X,Y),Z] + [B_R(Y,Z),X] + [B_R(Z,X),Y] = 0
  \tag{6.4.3}
\end{equation}
为使其成立，例如只须 $R$ 满足经典 Yang--Baxter 方程（classical Yang--Baxter equation）

% [原文 p.81]
\begin{equation}
  B_R(X,Y) = 0,\qquad \forall X,Y\in\mathfrak{g}
  \tag{6.4.4}
\end{equation}
即可。

\takeaway{R 矩阵给出新的 Lie 括号；其 Jacobi 恒等式归于经典 Yang--Baxter 方程。}

为理解经典 Yang--Baxter 方程，假定 $\mathfrak{g}$ 上存在内积
\begin{equation}
  (-,-):\mathfrak{g}\otimes\mathfrak{g}\to\R
  \tag{6.4.5}
\end{equation}
它关于 $\mathfrak{g}$ 的伴随作用不变，即
\begin{equation}
  ([X,Y],Z) + (X,[Y,Z]) = 0,\qquad \forall X,Y,Z\in\mathfrak{g}
  \tag{6.4.6}
\end{equation}
并取该内积下的规范正交基 $\{E_i\}_i$：
\begin{equation}
  (E_i,E_j) = \delta_{ij}.
  \tag{6.4.7}
\end{equation}
关于该基的 Lie 代数结构常数 $\{c_{ijk}\}_{i,j,k}$ 取为
\begin{equation}
  c_{ijk} := (E_i,[E_j,E_k])
  \tag{6.4.8}
\end{equation}
则由 Lie 括号的反称性，$c_{ijk}$ 关于 $j$ 与 $k$ 的交换是反称的：
\begin{equation}
  c_{ijk} = (E_i,[E_j,E_k]) = -(E_i,[E_k,E_j]) = -c_{ikj}.
  \tag{6.4.9}
\end{equation}
进一步，由内积 $(-,-)$ 的伴随不变性，$c_{ijk}$ 关于 $i$ 与 $k$ 的交换也是反称的：
\begin{equation}
  c_{ijk} = (E_i,[E_j,E_k]) = -([E_j,E_i],E_k) = -(E_k,[E_j,E_i]) = -c_{kji}.
  \tag{6.4.10}
\end{equation}

将 R 矩阵按分量写成
\begin{equation}
  R(E_j) = \sum_i r_{ij}E_i
  \tag{6.4.11}
\end{equation}
并定义 $r_{12},r_{13},r_{23},r_{32}\in U(\mathfrak{g})\otimes U(\mathfrak{g})\otimes U(\mathfrak{g})$ 为
\begin{align}
  r_{12} &= \sum_{i,j} r_{ij}E_i\otimes E_j\otimes 1, \tag{6.4.12}\\
  r_{13} &= \sum_{i,j} r_{ij}E_i\otimes 1\otimes E_j, \tag{6.4.13}\\
  r_{23} &= \sum_{i,j} r_{ij}1\otimes E_i\otimes E_j, \tag{6.4.14}\\
  r_{32} &= \sum_{i,j} r_{ij}1\otimes E_j\otimes E_i \tag{6.4.15}
\end{align}

% [原文 p.82]
\[
\begin{aligned}
R\bigl([RE_j,E_k]\bigr)
  &= R\Bigl(\Bigl[\sum_l r_{lj}E_l,\, E_k\Bigr]\Bigr) \\
  &= \sum_l r_{lj}R\bigl([E_l,E_k]\bigr) \\
  &= \sum_l r_{lj}R\Bigl(\sum_m c_{mlk}E_m\Bigr) \\
  &= \sum_{l,m} c_{mlk}r_{lj}R(E_m) \\
  &= \sum_{l,m} c_{mlk}r_{lj}\sum_i r_{im}E_i \\
  &= \sum_{l,m} c_{mlk}r_{im}r_{lj}E_i,
\end{aligned}
\]
% [Tutor Note: 末行照原书，求号只标 l,m 而 E_i 亦须对 i 求和；此处不改动。]

\[
\begin{aligned}
[r_{13},r_{32}]
  &= \Bigl[\sum_{i,m} r_{im}E_i\otimes 1\otimes E_m,\ \sum_{j,l} r_{lj}1\otimes E_j\otimes E_l\Bigr] \\
  &= \sum_{i,j,l,m} r_{im}r_{lj}\,E_i\otimes E_j\otimes [E_m,E_l] \\
  &= \sum_{i,j,k,l,m} r_{im}r_{lj}c_{kml}\,E_i\otimes E_j\otimes E_k
\end{aligned}
\]
成立，故由
\begin{equation}
  c_{kml} = -c_{mkl} = c_{mlk}
  \tag{6.4.16}
\end{equation}
注意可知
\begin{equation}
  \bigl(E_i, R([RE_j,E_k])\bigr) = \bigl(E_i\otimes E_j\otimes E_k,\ [r_{13},r_{32}]\bigr)
  \tag{6.4.17}
\end{equation}
同样地，
\begin{align}
  R([E_j,RE_k]) &= R\Bigl[E_j,\sum_l r_{lk}E_l\Bigr] \tag{6.4.18}\\
  &= \sum_{l,m} r_{lk}c_{mjl}R(E_m) \tag{6.4.19}\\
  &= \sum_{i,l,m} r_{im}r_{lk}c_{mjl}E_i, \tag{6.4.20}\\[2pt]
  [r_{12},r_{23}]
  &= \Bigl[\sum_{i,m} r_{im}E_i\otimes E_m\otimes 1,\ \sum_{l,k} r_{lk}1\otimes E_l\otimes E_k\Bigr] \tag{6.4.21}\\
  &= \sum_{i,j,k,l,m} r_{im}r_{lk}c_{jml}\,E_i\otimes E_j\otimes E_k \tag{6.4.22}
\end{align}
从而

% [原文 p.83]
\begin{equation}
  \bigl(E_i, R([E_j,RE_k])\bigr) = -\bigl(E_i\otimes E_j\otimes E_k,\ [r_{12},r_{23}]\bigr)
  \tag{6.4.23}
\end{equation}
并且
\begin{align}
  [R(E_j),R(E_k)] &= \Bigl[\sum_l r_{lj}E_l,\ \sum_m r_{mk}E_m\Bigr] \tag{6.4.24}\\
  &= \sum_{i,l,m} r_{lj}r_{mk}c_{ilm}E_i, \tag{6.4.25}\\[2pt]
  [r_{12},r_{13}]
  &= \Bigl[\sum_l r_{lj}E_l\otimes E_j\otimes 1,\ \sum_m r_{mk}E_m\otimes 1\otimes E_k\Bigr] \tag{6.4.26}\\
  &= \sum_{i,l,m} c_{ilm}r_{lj}r_{mk}\,E_i\otimes E_j\otimes E_k \tag{6.4.27}
\end{align}
从而
\begin{equation}
  \bigl(E_i,[R(E_j),R(E_k)]\bigr) = \bigl(E_i\otimes E_j\otimes E_k,\ [r_{12},r_{13}]\bigr)
  \tag{6.4.28}
\end{equation}
因此，经典 Yang--Baxter 方程被改写为
\begin{equation}
  [r_{12},r_{13}] - [r_{13},r_{32}] + [r_{12},r_{23}] = 0
  \tag{6.4.29}
\end{equation}
这里，若假定 R 矩阵的反称性
\begin{equation}
  (X,R(Y)) = -(Y,R(X))
  \tag{6.4.30}
\end{equation}
则
\begin{equation}
  r_{32} = -r_{23}
  \tag{6.4.31}
\end{equation}
成立，经典 Yang--Baxter 方程即为
\begin{equation}
  [r_{12},r_{13}] + [r_{12},r_{23}] + [r_{13},r_{23}] = 0
  \tag{6.4.32}
\end{equation}
所给出。

另一方面，当 $A$ 是量子群（或更一般地，拟三角 Hopf 代数）时，存在 $R\in A\otimes A$ 满足 $A\otimes A\otimes A$ 中的 Yang--Baxter 方程
\begin{equation}
  R_{12}R_{13}R_{23} = R_{23}R_{13}R_{12}
  \tag{6.4.33}
\end{equation}
当 $A$ 定义在形式幂级数环 $\C[[\hbar]]$ 上、且 $R$ 具有
\begin{equation}
  R = 1 + \hbar r + \hbar^2 s + o(\hbar^3)
  \tag{6.4.34}
\end{equation}
的形状时，忽略 $\hbar$ 的 3 次以上项，则

% [原文 p.84]
\[
\begin{aligned}
R_{12}R_{13}R_{23}
  &= (1+\hbar r_{12}+\hbar^2 s_{12})(1+\hbar r_{13}+\hbar^2 s_{13})(1+\hbar r_{23}+\hbar^2 s_{23}) \\
  &= 1 + \hbar(r_{12}+r_{13}+r_{23}) \\
  &\qquad + \hbar^2\bigl(r_{12}r_{13} + r_{12}r_{23} + r_{13}r_{23} + s_{12} + s_{13} + s_{23}\bigr), \\
R_{23}R_{13}R_{12}
  &= 1 + \hbar(r_{23}+r_{13}+r_{12}) \\
  &\qquad + \hbar^2\bigl(r_{23}r_{13} + r_{23}r_{12} + r_{13}r_{12} + s_{23} + s_{13} + s_{12}\bigr)
\end{aligned}
\]
从而
\begin{equation}
  r_{12}r_{13} + r_{12}r_{23} + r_{13}r_{23}
  = r_{23}r_{13} + r_{23}r_{12} + r_{13}r_{12}
  \tag{6.4.35}
\end{equation}
成立，而这正是经典 Yang--Baxter 方程本身。
% [Tutor Note: 量子 Yang--Baxter 方程 (6.4.33) 的 $\hbar$ 一次项给出经典方程，二次项给出 $s$ 的相容条件；此处只取一次项。]

\section{AKS 定理}

当 Lie 代数 $\mathfrak{g}$ 上给定 R 矩阵 $R\in\End(\mathfrak{g})$ 时，用 $\mathfrak{g}_R$ 表示在向量空间 $\mathfrak{g}$ 上按 (6.4.1) 赋予 Lie 代数结构所得的 Lie 代数。给定 Lie 代数 $\mathfrak{g}$ 的 R 矩阵 $R$ 后，对偶空间 $\mathfrak{g}^*$ 上便确定两个 Lie--Poisson 结构 $\{-,-\}$ 与 $\{-,-\}_R$。给定泊松流形 $(M,\{-,-\})$ 时，$M$ 上与任意函数都泊松可交换的函数称为 \textbf{Casimir 函数}（Casimir function）。Lie--Poisson 流形 $(\mathfrak{g}^*,\{-,-\})$ 上 Casimir 函数全体构成的环记作 $I(\mathfrak{g})$。

\takeaway{Casimir 函数就是余伴随不变函数；AKS 定理把它的 Hamilton 流写成 Lax 方程。}

\begin{proposition*}{6.5.1}
  $\mathfrak{g}^*$ 上的函数 $\varphi$ 属于 $I(\mathfrak{g})$ 的充要条件是：$\varphi$ 关于 $G$ 的余伴随作用不变。
\end{proposition*}

\begin{proof}
  由定义，$\mathfrak{g}^*$ 上的函数 $\varphi$ 属于 $I(\mathfrak{g})$ 的充要条件是：对任意 $H\in C^\infty(\mathfrak{g}^*)$ 与 $\xi\in\mathfrak{g}^*$，
  \begin{equation}
    \{H,\varphi\}(\xi) = \bigl([(dH)_\xi,(d\varphi)_\xi],\ \xi\bigr) = 0
    \tag{6.5.1}
  \end{equation}
  成立。借助自然同构 $T^*_\xi\mathfrak{g}^*\cong\mathfrak{g}$，把 $(dH)_\xi\in T^*_\xi\mathfrak{g}^*$ 视为 $\mathfrak{g}$ 的元素并记为 $X$，则该条件可改写为
  \begin{equation}
    \bigl\langle [X,(d\varphi)_\xi],\ \xi\bigr\rangle = \bigl\langle \ad_X(d\varphi)_\xi,\ \xi\bigr\rangle = 0
    \tag{6.5.2}
  \end{equation}
  任意 $X\in\mathfrak{g}$ 都可由适当的 $H\in C^\infty(\mathfrak{g}^*)$ 表为 $X=(dH)_\xi$，故 (6.5.2) 必须对任意 $X\in\mathfrak{g}$ 成立，而这意味着 $\varphi$ 关于 $G$ 的余伴随作用不变。
\end{proof}
% [Tutor Note: (6.5.1) 即 Lie--Poisson 括号在 $\mathfrak{g}^*$ 上的定义式 $\{H,\varphi\}(\xi)=\langle\xi,[(dH)_\xi,(d\varphi)_\xi]\rangle$。]

(6.5.2) 对任意 $\xi\in\mathfrak{g}^*$ 与 $X\in\mathfrak{g}$ 而言，等价于

% [原文 p.85]
\begin{equation}
  \bigl\langle X,\ \ad^*_{(d\varphi)_\xi}\xi\bigr\rangle = 0,
  \tag{6.5.3}
\end{equation}
亦即对任意 $\xi\in\mathfrak{g}^*$ 都有
\begin{equation}
  \ad^*_{(d\varphi)_\xi}\xi = 0
  \tag{6.5.4}
\end{equation}
也可这样表述。设 $I(\mathfrak{g})$ 中函数无关的元的个数为 $l$。当 $\mathfrak{g}$ 为半单 Lie 代数时，$l$ 与 $\mathfrak{g}$ 的秩一致。

下述定理把发现者名字的首字母并列，称为 \textbf{AKS 定理}（AKS theorem）：
% [Tutor Note: AKS 取自 Adler、Kostant、Symes 三人姓氏的首字母，三人在 1970 年代各自独立得到该结论。]

\begin{theorem*}{6.5.2}
  $\varphi\in I(\mathfrak{g})$ 关于 $\{-,-\}_R$ 的 Hamilton 流，对 $L\in\mathfrak{g}^*$ 由
  \begin{equation}
    \frac{\rd L}{\rd t} = -\ad^*_M L,\qquad
    M = \frac{1}{2}R(d\varphi)_L
    \tag{6.5.5}
  \end{equation}
  给出。
\end{theorem*}

\begin{proof}
  关于 $\psi\in C^\infty(\mathfrak{g}^*)$ 的运动方程为
  \begin{align}
    \frac{\rd\psi}{\rd t}(L) &= \{\varphi,\psi\}_R(L) \tag{6.5.6}\\
    &= \bigl\langle L,\ [(d\varphi)_L,(d\psi)_L]_R\bigr\rangle \tag{6.5.7}\\
    &= \Bigl\langle L,\ \frac{1}{2}\bigl([R(d\varphi)_L,(d\psi)_L] + [(d\varphi)_L,R(d\psi)_L]\bigr)\Bigr\rangle \tag{6.5.8}\\
    &= \Bigl\langle L,\ \frac{1}{2}[R(d\varphi)_L,(d\psi)_L]\Bigr\rangle \tag{6.5.9}\\
    &= \frac{1}{2}\bigl\langle L,\ \ad_{R(d\varphi)_L}(d\psi)_L\bigr\rangle \tag{6.5.10}\\
    &= -\frac{1}{2}\bigl\langle \ad^*_{R(d\varphi)_L}L,\ (d\psi)_L\bigr\rangle \tag{6.5.11}
  \end{align}
  其中，由 (6.5.8) 到 (6.5.9) 的变形用到了来自 $\varphi\in I(\mathfrak{g})$ 的 (6.5.2)。
\end{proof}

下述系称为对合性定理（involutivity theorem）：
% [Tutor Note: 原书作「包含性定理」，西文原注 involutivity；通行译名为「对合性」，此处从通行译法。]

\begin{corollary*}{6.5.3}
  对任意 $\varphi,\psi\in I(\mathfrak{g})$ 都有 $\{\varphi,\psi\}_R = 0$。
\end{corollary*}

AKS 定理主张：R 矩阵型的 Lie--Poisson 流形上，由 Casimir 函数给出的 Hamilton 流具有 Lax 形式的写法。

\begin{corollary*}{6.5.4}
  当 $\mathfrak{g}$ 具有非退化不变内积时，以该内积把 $\mathfrak{g}$ 与 $\mathfrak{g}^*$ 视为同一，则取 $\varphi\in I(\mathfrak{g})$ 为 Hamilton 函数的 Hamilton 系具有 Lax 表示
% [原文 p.86]
  \begin{equation}
    \frac{\rd L}{\rd t} = [L,M]
    \tag{6.5.12}
  \end{equation}
\end{corollary*}

因此，若 $\mathfrak{g}^*$ 的 $\ad^*_{\mathfrak{g}_R}$ 轨道 $\mathcal{O}$ 是 $2l$ 维的，并且把 $I(\mathfrak{g})$ 中函数无关的元限制到 $\mathcal{O}$ 上仍然函数无关，则对 $I(\mathfrak{g})$ 的任意元 $\varphi$，以 $\varphi$ 为 Hamilton 函数、定义在 $\mathcal{O}$ 上的 Hamilton 系便是具有 Lax 表示的 Liouville 可积系统。

\section{因子化定理}

考虑 Lie 代数 $\mathfrak{g}$ 作为向量空间是两部分子 Lie 代数 $\mathfrak{g}_+$ 与 $\mathfrak{g}_-$ 的直和的情形：
\begin{equation}
  \mathfrak{g} = \mathfrak{g}_+ \oplus \mathfrak{g}_-.
  \tag{6.6.1}
\end{equation}
这是作为向量空间的直和分解，并非作为 Lie 代数的直和分解，故虽有 $[\mathfrak{g}_+,\mathfrak{g}_+]\subset\mathfrak{g}_+$ 且 $[\mathfrak{g}_-,\mathfrak{g}_-]\subset\mathfrak{g}_-$，但 $[\mathfrak{g}_+,\mathfrak{g}_-]=0$ 未必成立。设 $P_\pm:\mathfrak{g}\to\mathfrak{g}_\pm$ 为直和分解 (6.6.1) 所附带的投影算子，令 $R=P_+-P_-$，则由 (6.4.1)
\begin{align}
  [X,Y]_R &= \frac{1}{2}\bigl([RX,Y] + [X,RY]\bigr) \tag{6.6.2}\\
  &= \frac{1}{2}\bigl([X_+-X_-,\ Y_++Y_-] + [X_++X_-,\ Y_+-Y_-]\bigr) \tag{6.6.3}\\
  &= [X_+,Y_+] - [X_-,Y_-] \tag{6.6.4}
\end{align}
成立。这里
\begin{equation}
  X_\pm = P_\pm X,\qquad Y_\pm = P_\pm Y
  \tag{6.6.5}
\end{equation}
而 (6.6.4) 显然在 $\mathfrak{g}$ 上确定一个 Lie 代数结构，故 $R$ 是 R 矩阵。另一方面，对应于 Lie 代数 $\mathfrak{g}$ 的子 Lie 代数 $\mathfrak{g}_+$ 与 $\mathfrak{g}_-$，Lie 群 $G$ 的子 Lie 群 $G_+$ 与 $G_-$ 也随之确定。

\takeaway{$\mathfrak{g}=\mathfrak{g}_+\oplus\mathfrak{g}_-$ 给出 R 矩阵；Casimir 流由群上的因子化 $e^{tX_0}=g_+^{-1}g_-$ 显式解出。}

\begin{theorem*}{6.6.1}
  选定并固定 $\varphi\in I(\mathfrak{g})$ 与 $L_0\in\mathfrak{g}^*$，令 $X_0=(d\varphi)_{L_0}\in\mathfrak{g}$。再假定存在 $g_+:\R\to G_+$ 与 $g_-:\R\to G_-$，使得对任意 $t\in\R$ 都有
  \begin{equation}
    \exp(tX_0) = g_+(t)^{-1}g_-(t)
    \tag{6.6.6}
  \end{equation}
  成立。此时，以 $L_0$ 为初始值、以 $\varphi$ 为 Hamilton 函数的 Hamilton 流由
  \begin{equation}
    L(t) = \operatorname{Ad}^*_{g_+(t)}L_0
    \tag{6.6.7}
  \end{equation}
  给出。
\end{theorem*}

% [原文 p.87]
\begin{proof}
  一般地，给定 $g:\R\to G$ 时，对
  \begin{equation}
    gg^{-1} = e
    \tag{6.6.8}
  \end{equation}
  的两边关于 $t$ 求导，得
  \begin{equation}
    \dot{g}g^{-1} + g\frac{\rd(g^{-1})}{\rd t} = 0
    \tag{6.6.9}
  \end{equation}
  即
  \begin{equation}
    \frac{\rd(g^{-1})}{\rd t} = -g^{-1}\dot{g}g^{-1}
    \tag{6.6.10}
  \end{equation}
  对常数 $X_0\in\mathfrak{g}$，令
  \begin{equation}
    X = \operatorname{Ad}_{g^{-1}}X_0 = g^{-1}X_0g
    \tag{6.6.11}
  \end{equation}
  则
  \begin{align}
    \dot{X} &= \frac{\rd g^{-1}}{\rd t}X_0g + g^{-1}X_0\dot{g} \tag{6.6.12}\\
    &= -g^{-1}\dot{g}g^{-1}X_0g + g^{-1}X_0\dot{g} \tag{6.6.13}\\
    &= -g^{-1}\bigl(\dot{g}g^{-1}X_0 - X_0\dot{g}g^{-1}\bigr)g \tag{6.6.14}\\
    &= -\operatorname{Ad}_{g^{-1}}\ad_{\dot{g}g^{-1}}X_0 \tag{6.6.15}
  \end{align}
  成立。这里须注意 $\dot{g}\in T_gG$，从而 $\dot{g}g^{-1}\in T_eG=\mathfrak{g}$。其次，对余伴随表示的定义式
  \begin{equation}
    \bigl\langle \operatorname{Ad}^*_g\xi_0,\ X_0\bigr\rangle
    = \bigl\langle \xi_0,\ \operatorname{Ad}_{g^{-1}}X_0\bigr\rangle
    \tag{6.6.16}
  \end{equation}
  的两边关于 $t$ 求导，得
  \begin{align}
    \Bigl\langle \frac{\rd}{\rd t}\operatorname{Ad}^*_g\xi_0,\ X_0\Bigr\rangle
      &= \bigl\langle \xi_0,\ -\operatorname{Ad}_{g^{-1}}\ad_{\dot{g}g^{-1}}X_0\bigr\rangle \tag{6.6.17}\\
    &= \bigl\langle \operatorname{Ad}^*_g\xi_0,\ -\ad_{\dot{g}g^{-1}}X_0\bigr\rangle \tag{6.6.18}\\
    &= \bigl\langle \ad_{\dot{g}g^{-1}}\operatorname{Ad}^*_g\xi_0,\ X_0\bigr\rangle \tag{6.6.19}
  \end{align}
  也就是说，对常数 $\xi_0\in\mathfrak{g}^*$，令
  \begin{equation}
    \xi = \operatorname{Ad}^*_g\xi_0
    \tag{6.6.20}
  \end{equation}
  则
  \begin{equation}
    \dot{\xi} = \ad_{\dot{g}g^{-1}}\xi
    \tag{6.6.21}
  \end{equation}
  成立。因此，对 (6.6.7) 求导得
  \begin{equation}
    \dot{L} = \ad^*_{\dot{g}_+g_+^{-1}}L
    \tag{6.6.22}
  \end{equation}

% [原文 p.88]
一般地，当 Lie 群 $G$ 作用于流形 $M$、且 $M$ 上的函数 $f$ 在 $G$ 的作用下不变时，对任意 $X\in T_xM$ 与 $g\in G$ 都有
\begin{equation}
  (df)_{gx}(g_*X) = (df)_x(X)
  \tag{6.6.23}
\end{equation}
成立。换言之，
\begin{equation}
  (df)_{gx} = (g^{-1})^*(df)_x
  \tag{6.6.24}
\end{equation}
成立。注意 $\operatorname{Ad}^*_g:\mathfrak{g}^*\to\mathfrak{g}^*$ 是 $\operatorname{Ad}_{g^{-1}}:\mathfrak{g}\to\mathfrak{g}$ 的转置，把 (6.6.24) 应用于 $\mathfrak{g}^*$ 上 $\operatorname{Ad}^*_G$ 不变函数 $\varphi$，得
\begin{align}
  (d\varphi)_{L(t)} &= (d\varphi)_{\operatorname{Ad}^*_{g_+}L_0} \tag{6.6.25}\\
  &= \operatorname{Ad}_{g_+}(d\varphi)_{L_0} \tag{6.6.26}\\
  &= \operatorname{Ad}_{g_+}X_0 \tag{6.6.27}
\end{align}
另一方面，对 (6.6.6) 的两边关于 $t$ 求导，得
\begin{align}
  X_0\exp(tX_0) &= \frac{\rd g_+^{-1}}{\rd t}g_- + g_+^{-1}\frac{\rd g_-}{\rd t} \tag{6.6.28}\\
  &= -g_+^{-1}\dot{g}_+g_+^{-1}g_- + g_+^{-1}\dot{g}_- \tag{6.6.29}
\end{align}
因此
\begin{equation}
  X_0g_+^{-1}g_- = -g_+^{-1}\dot{g}_+g_+^{-1}g_- + g_+^{-1}\dot{g}_-
  \tag{6.6.30}
\end{equation}
成立，在此式两边左乘 $g_+$、右乘 $g_-^{-1}$，得
\begin{equation}
  g_+X_0g_+^{-1} = -\dot{g}_+g_+^{-1} + \dot{g}_-g_-^{-1}
  \tag{6.6.31}
\end{equation}
即
\begin{equation}
  \operatorname{Ad}_{g_+}X_0 = -\dot{g}_+g_+^{-1} + \dot{g}_-g_-^{-1}
  \tag{6.6.32}
\end{equation}
这里，
\begin{equation}
  \dot{g}_+g_+^{-1}\in\mathfrak{g}_+,\qquad \dot{g}_-g_-^{-1}\in\mathfrak{g}_-
  \tag{6.6.33}
\end{equation}
故
\begin{equation}
  \bigl(\operatorname{Ad}_{g_+}X_0\bigr)_+ = -\dot{g}_+g_+^{-1}
  \tag{6.6.34}
\end{equation}
成立。由 (6.6.27) 与 (6.6.34) 得
\begin{equation}
  \bigl((d\varphi)_L\bigr)_+ = -\dot{g}_+g_+^{-1}
  \tag{6.6.35}
\end{equation}
另一方面，由于 $\varphi$ 是 Casimir 函数，由 (6.5.4) 得

% [原文 p.89]
\begin{align}
  0 &= \ad^*_{(d\varphi)_\xi}\xi \tag{6.6.36}\\
  &= \ad^*_{((d\varphi)_\xi)_+ + ((d\varphi)_\xi)_-}\xi \tag{6.6.37}\\
  &= \ad^*_{((d\varphi)_\xi)_+}\xi + \ad^*_{((d\varphi)_\xi)_-}\xi \tag{6.6.38}
\end{align}
从而
\begin{equation}
  \ad^*_{((d\varphi)_\xi)_-}\xi = -\ad^*_{((d\varphi)_\xi)_+}\xi
  \tag{6.6.39}
\end{equation}
利用这一点，
\begin{align}
  -\ad^*_ML &= -\ad^*_{\frac{1}{2}R(d\varphi)_L}L \tag{6.6.40}\\
  &= -\ad^*_{\frac{1}{2}(((d\varphi)_L)_+ - ((d\varphi)_L)_-)}L \tag{6.6.41}\\
  &= -\ad^*_{((d\varphi)_L)_+}L \tag{6.6.42}\\
  &= \ad^*_{\dot{g}_+g_+^{-1}}L \tag{6.6.43}
\end{align}
  成立。这里，第一个等号用了 $M$ 的定义，下一个等号用了 $R$ 的定义，再下一个等号用了 (6.6.39)，最后一个等号用了 (6.6.35)。将此式与 (6.6.22) 比较，定理 6.6.1 便得证。
\end{proof}

\section{Toda 格子再论}

设 Lie 代数 $\mathfrak{g}$ 具有作为向量空间的直和分解
\begin{equation}
  \mathfrak{g} = \bigoplus_{i\in\Z}\mathfrak{g}_i
  \tag{6.7.1}
\end{equation}
并满足
\begin{equation}
  [\mathfrak{g}_i,\mathfrak{g}_j] \subset \mathfrak{g}_{i+j}
  \tag{6.7.2}
\end{equation}
此时称 $\mathfrak{g}$ 为 \textbf{分次 Lie 代数}（graded Lie algebra）。对分次 Lie 代数 $\mathfrak{g}$，令
\begin{equation}
  \mathfrak{g}_+ := \bigoplus_{i=0}^{\infty}\mathfrak{g}_i,\qquad
  \mathfrak{g}_- := \bigoplus_{i=1}^{\infty}\mathfrak{g}_{-i}
  \tag{6.7.3}
\end{equation}
则 $\mathfrak{g}_+$ 与 $\mathfrak{g}_-$ 都是 $\mathfrak{g}$ 的子 Lie 代数。将直和分解 $\mathfrak{g}=\mathfrak{g}_+\oplus\mathfrak{g}_-$ 所附带的投影写作
\begin{equation}
  P_\pm:\mathfrak{g}\to\mathfrak{g}_\pm
  \tag{6.7.4}
\end{equation}
并以
\begin{equation}
  R = P_+ - P_-
  \tag{6.7.5}
\end{equation}
定义直和分解型的 R 矩阵。Lie 代数 $\mathfrak{g}_R$ 的余伴随作用记作 $\ad^*_R$。

% [原文 p.90]
\takeaway{分次结构给出直和分解型 R 矩阵；其轨道上的 Toda 流化为 Lax 方程。}

\begin{lemma*}{6.7.1}
  对任意 $X\in\mathfrak{g}$ 与 $\xi\in\mathfrak{g}^*$ 都有
  \begin{equation}
    \ad^*_{RX}\xi = \bigl(\ad^*_{X_+}\xi\bigr)_+ - \bigl(\ad^*_{X_-}\xi\bigr)_-
    \tag{6.7.6}
  \end{equation}
  成立。
\end{lemma*}

\begin{proof}
  按直和分解 $\mathfrak{g}^*=\mathfrak{g}^*_+\oplus\mathfrak{g}^*_-$ 把 $\xi\in\mathfrak{g}^*$ 写成
  \begin{equation}
    \xi = \xi_+ + \xi_-,\qquad \xi_\pm\in\mathfrak{g}^*_\pm
    \tag{6.7.7}
  \end{equation}
  则
  \begin{equation}
    \ad_{X_+}Y_+ = [X_+,Y_+]\in\mathfrak{g}_+
    \tag{6.7.8}
  \end{equation}
  故
  \begin{equation}
    \bigl\langle \ad^*_{X_+}\xi_-,\ Y_+\bigr\rangle
    = -\bigl\langle \xi_-,\ \ad_{X_+}Y_+\bigr\rangle
    = 0,
    \tag{6.7.9}
  \end{equation}
  即
  \begin{equation}
    \bigl(\ad^*_{X_+}\xi_-\bigr)_+ = 0
    \tag{6.7.10}
  \end{equation}
  因此
  \begin{align}
    \bigl(\ad^*_{X_+}\xi\bigr)_+
      &= \bigl(\ad^*_{X_+}\xi_+\bigr)_+ + \bigl(\ad^*_{X_+}\xi_-\bigr)_+ \tag{6.7.11}\\
      &= \bigl(\ad^*_{X_+}\xi_+\bigr)_+ \tag{6.7.12}
  \end{align}
  同样，
  \begin{equation}
    \bigl(\ad^*_{X_-}\xi\bigr)_- = \bigl(\ad^*_{X_-}\xi_-\bigr)_-
    \tag{6.7.13}
  \end{equation}
  因此
  \[
  \begin{aligned}
    \bigl(\ad^*_{X_+}\xi\bigr)_+ - \bigl(\ad^*_{X_-}\xi\bigr)_-
      &= \bigl(\ad^*_{X_+}\xi_+\bigr)_+ - \bigl(\ad^*_{X_-}\xi_-\bigr)_- \\
      &= \ad^*_{RX}\xi
  \end{aligned}
  \]
  成立。
\end{proof}

于是，对由分次 Lie 代数得到的直和分解型 R 矩阵，下述结论成立：

\begin{corollary*}{6.7.2}
  对非负整数 $m,n$ 与 $f\in\mathfrak{g}^*_{-m-1}$，$\mathfrak{g}^*$ 的仿射子空间
  \begin{equation}
    f + \bigoplus_{i=-m}^{n-1}\mathfrak{g}^*_i
    \tag{6.7.14}
  \end{equation}
  是泊松子空间。
\end{corollary*}

% [原文 p.91]
\begin{proof}
  对任意 $X\in\mathfrak{g}_i$、$\xi\in\mathfrak{g}^*_j$ 以及 $Y\in\mathfrak{g}_k$，
  \begin{equation}
    \ad_XY = [X,Y]\in\mathfrak{g}_{i+k}
    \tag{6.7.15}
  \end{equation}
  故
  \begin{equation}
    \bigl\langle \ad^*_X\xi,\ Y\bigr\rangle = -\bigl\langle \xi,\ \ad_XY\bigr\rangle
    \tag{6.7.16}
  \end{equation}
  只在 $i+k=j$ 时才可能取非零值。也就是说，
  \begin{equation}
    \ad^*_{\mathfrak{g}_i}\mathfrak{g}^*_j \subset \mathfrak{g}^*_{j-i}
    \tag{6.7.17}
  \end{equation}
  由 Lie--Poisson 括号的定义，Hamilton 向量场构成的切分布由余伴随作用的基本向量场生成，故由引理 3.6.4，只须证明 $\mathfrak{g}^*$ 的仿射子空间 (6.7.14) 关于 $\ad^*_R$ 不变。对任意 $X=X_++X_-\in\mathfrak{g}_+\oplus\mathfrak{g}_-$ 与 $\xi\in f+\bigoplus_{i=-m}^{n-1}\mathfrak{g}^*_i$，由 (6.7.17)
  \begin{equation}
    \ad^*_{X_+}\xi\in\mathfrak{g}^*_{\le n-1},\qquad
    \ad^*_{X_-}\xi\in\mathfrak{g}^*_{\ge -m}
    \tag{6.7.18}
  \end{equation}
  因此
  \begin{equation}
    \bigl(\ad^*_{X_+}\xi\bigr)_+ \in \bigoplus_{i=0}^{m-1}\mathfrak{g}^*_i,
    \tag{6.7.19}
  \end{equation}
  % [Tutor Note: (6.7.19) 原书右端上界印作 $m-1$；按 (6.7.18) 与 $\xi\in f+\bigoplus_{i=-m}^{n-1}\mathfrak{g}^*_i$ 推来应为 $n-1$，此处照录原文。]
  \begin{equation}
    \bigl(\ad^*_{X_-}\xi\bigr)_- \in \bigoplus_{i=-m}^{-1}\mathfrak{g}^*_i
    \tag{6.7.20}
  \end{equation}
  并且由 (6.7.6)
  \begin{equation}
    \ad^*_{RX}\xi \in \bigoplus_{i=-m}^{n-1}\mathfrak{g}^*_i
    \tag{6.7.21}
  \end{equation}
  成立。
\end{proof}

取 $\mathfrak{g}=\mathfrak{sl}(n,\R)$，考虑直和分解
\begin{equation}
  \mathfrak{g} = \bigoplus_{i=-n+1}^{n-1}\mathfrak{g}_i,\qquad
  \mathfrak{g}_i = \Bigl\{\ \sum_{k,l=1}^{n}X_{kl}E_{kl}\ \Bigm|\
  l-k\ne i\ \text{时}\ X_{kl}=0\ \Bigr\}
  \tag{6.7.22}
\end{equation}
也就是说，$\mathfrak{g}_i$ 是由与对角线相距 $i$ 的分量以外全为零的矩阵构成的。按不变内积
\begin{equation}
  (X,Y) = \tr(XY) = \sum_{i,j}X_{ij}Y_{ji}
  \tag{6.7.23}
\end{equation}
把 $\mathfrak{g}^*_i$ 与 $\mathfrak{g}_{-i}$ 视为同一。在系 6.7.2 中取 $m=0$、$n=2$，并令 $f$ 为

% [原文 p.92]
\begin{equation}
  f = \begin{pmatrix}
    0&1&0&\cdots&0\\
    0&0&1&\cdots&0\\
    \vdots&\vdots&\vdots&\ddots&\vdots\\
    0&0&0&\cdots&1\\
    0&0&0&\cdots&0
  \end{pmatrix}
  \in\mathfrak{g}_1\cong\mathfrak{g}^*_{-1}
  \tag{6.7.24}
\end{equation}
则子空间 (6.7.14) 为
\begin{equation}
  \left\{
  \begin{pmatrix}
    b_1&1&0&\cdots&0&0\\
    a_1&b_2&1&\cdots&0&0\\
    0&a_2&b_3&\cdots&0&0\\
    \vdots&\vdots&\ddots&\ddots&\vdots&\vdots\\
    0&0&\cdots&a_{n-2}&b_{n-1}&1\\
    0&0&\cdots&0&a_{n-1}&b_n
  \end{pmatrix}
  \ \middle|\ \sum_{i=1}^{n}b_i = 0
  \right\}
  \tag{6.7.25}
\end{equation}
这里，条件 $\sum_{i=1}^{n}b_i=0$ 来自
\begin{equation}
  \mathfrak{sl}(n,\R) = \{X\in\mathrm{Mat}_{n\times n}(\R)\mid \tr X = 0\}
  \tag{6.7.26}
\end{equation}

\begin{proposition*}{6.7.3}
  子空间 (6.7.25) 上坐标函数 $a_i$ 与 $b_j$ 的 Poisson 括号为
  \begin{equation}
    \{a_i,a_j\}_R = \{b_i,b_j\}_R = 0,\qquad
    \{a_i,b_j\}_R = (\delta_{i+1,j} - \delta_{ij})a_i
    \tag{6.7.27}
  \end{equation}
\end{proposition*}

\begin{proof}
  记子空间 (6.7.25) 为 $M$，$M$ 上的点写作 $\xi=(a,b)=(a_1,\dots,a_{n-1},b_1,\dots,b_n)$。坐标函数 $a_i\in C^\infty(M)$ 在点 $\xi=(a,b)\in M$ 处的微分 $\rd a_i$ 经自然嵌入 $T^*_\xi M\subset\mathfrak{g}$ 对应于 $E_{i,i+1}\in\mathfrak{g}_1$。由 Lie--Poisson 括号的定义，
  \begin{align}
    \{a_i,a_j\}_R(\xi) &= \bigl\langle [\rd a_i,\rd a_j]_R,\ \xi\bigr\rangle \tag{6.7.28}\\
    &= \bigl\langle [E_{i,i+1},E_{j,j+1}]_R,\ \xi\bigr\rangle \tag{6.7.29}\\
    &= \bigl\langle (E_{i,i+1}E_{j,j+1} - E_{j,j+1}E_{i,i+1})_+,\ \xi\bigr\rangle \tag{6.7.30}\\
    &= \bigl\langle (\delta_{i+1,j}E_{i,j+1} - \delta_{j+1,i}E_{j,i+1})_+,\ \xi\bigr\rangle \tag{6.7.31}\\
    &= \bigl\langle \delta_{i+1,j}E_{i,i+2} - \delta_{j+1,i}E_{j,j+2},\ \xi\bigr\rangle \tag{6.7.32}\\
    &= 0 \tag{6.7.33}
  \end{align}
  成立。其中最后的等号来自：对任意 $i$，$\xi$ 的 $(i+2,i)$ 分量为零。同样地，
  \begin{equation}
    [E_{i,i},E_{j,j}] = 0
    \tag{6.7.34}
  \end{equation}

% [原文 p.93]
从而
\begin{equation}
  \{b_i,b_j\}_R(\xi) = \bigl\langle [\rd b_i,\rd b_j]_R,\ \xi\bigr\rangle
  = \bigl\langle [E_{ii},E_{jj}]_R,\ \xi\bigr\rangle = 0
  \tag{6.7.35}
\end{equation}
并且
\begin{align}
  [E_{i,i+1},E_{jj}] &= E_{i,i+1}E_{jj} - E_{jj}E_{i,i+1} \tag{6.7.36}\\
  &= \delta_{i+1,j}E_{ij} - \delta_{ji}E_{j,i+1} \tag{6.7.37}\\
  &= \delta_{i+1,j}E_{i,i+1} - \delta_{ij}E_{i,i+1} \tag{6.7.38}\\
  &= (\delta_{i+1,j} - \delta_{ij})E_{i,i+1} \tag{6.7.39}
\end{align}
故
\begin{align}
  \{a_i,b_j\}_R(\xi) &= \bigl\langle [\rd a_i,\rd b_j]_R,\ \xi\bigr\rangle \tag{6.7.40}\\
  &= \bigl\langle [E_{i,i+1},E_{jj}]_R,\ \xi\bigr\rangle \tag{6.7.41}\\
  &= \bigl\langle \bigl((\delta_{i+1,j}-\delta_{ij})E_{i,i+1}\bigr)_+,\ \xi\bigr\rangle \tag{6.7.42}\\
  &= (\delta_{i+1,j}-\delta_{ij})a_i \tag{6.7.43}
\end{align}
成立。
\end{proof}

令
\begin{equation}
  e = \begin{pmatrix}
    0&0&0&\cdots&0\\
    1&0&0&\cdots&0\\
    0&1&0&\cdots&0\\
    \vdots&\vdots&\ddots&\ddots&\vdots\\
    0&0&\cdots&1&0
  \end{pmatrix}
  \tag{6.7.44}
\end{equation}
并把 $e+f$ 在 (6.7.25) 中的 $\ad^*_R$ 轨道记作 $\mathcal{O}$，则
\[
\mathcal{O} = \Bigl\{\ \xi=(a,b)\ \Bigm|\
\text{对任意 } i=1,\dots,n-1 \text{ 有 } a_i>0
\text{ 且 } \sum_{i=1}^{n}b_i = 0\ \Bigr\}
\]
由命题 6.7.3，可取 $\mathcal{O}$ 的 Darboux 坐标为
\begin{equation}
  a_i = e^{q_i-q_{i+1}},\qquad b_i = p_i
  \tag{6.7.45}
\end{equation}
其中
\begin{equation}
  \sum_{i=1}^{n}p_i = \sum_{i=1}^{n}q_i = 0
  \tag{6.7.46}
\end{equation}
取 Casimir 函数
\begin{equation}
  H(X) = \frac{1}{2}\tr(X^2)
  \tag{6.7.47}
\end{equation}
为 Hamilton 函数，则它在 Darboux 坐标下

% [原文 p.94]
\begin{equation}
  H = \frac{1}{2}\sum_{i=1}^{n}p_i^2 + \sum_{i=1}^{n-1}e^{q_i-q_{i+1}}
  \tag{6.7.48}
\end{equation}
可以表出。Lax 方程 (6.5.12) 借助
\begin{equation}
  M_\pm = \frac{1}{2}(R+1)(d\varphi)_L
  \tag{6.7.49}
\end{equation}
可改写为
\begin{equation}
  \frac{\rd L}{\rd t} = [L,M_\pm]
  \tag{6.7.50}
\end{equation}
此时的 Lax 对为
\begin{equation}
  L = \begin{pmatrix}
    p_1&1&0&\cdots&0&0\\
    e^{q_1-q_2}&p_2&1&\cdots&0&0\\
    0&e^{q_2-q_3}&p_3&\ddots&0&0\\
    \vdots&\vdots&\ddots&\ddots&\ddots&\vdots\\
    0&0&\cdots&e^{q_{n-2}-q_{n-1}}&p_{n-1}&1\\
    0&0&\cdots&0&e^{q_{n-1}-q_n}&p_n
  \end{pmatrix},
  \tag{6.7.51}
\end{equation}
\begin{equation}
  M_+ = \begin{pmatrix}
    p_1&1&0&\cdots&0&0\\
    0&p_2&1&\cdots&0&0\\
    0&0&p_3&\cdots&0&0\\
    \vdots&\vdots&\vdots&\ddots&\vdots&\vdots\\
    0&0&0&\cdots&p_{n-1}&1\\
    0&0&0&\cdots&0&p_n
  \end{pmatrix},
  \tag{6.7.52}
\end{equation}
\begin{equation}
  M_- = -\begin{pmatrix}
    0&0&\cdots&0&0\\
    e^{q_1-q_2}&0&\cdots&0&0\\
    0&e^{q_2-q_3}&\cdots&0&0\\
    \vdots&\vdots&\ddots&\vdots&\vdots\\
    0&0&\cdots&e^{q_{n-1}-q_n}&0
  \end{pmatrix},
  \tag{6.7.53}
\end{equation}
% [Tutor Note: (6.7.49) 原书右端对 $M_+$、$M_-$ 只印一个式子 $\frac{1}{2}(R+1)(d\varphi)_L$，结合 (6.7.52)(6.7.53) 应理解为该式的 $\pm$ 分量；此处照录原文。]

Toda 格子对一般的单 Lie 代数也可定义。设 $\mathfrak{g}$ 的根空间分解为
\begin{equation}
  \mathfrak{g} = \mathfrak{a} \oplus \bigoplus_{\alpha\in\Delta}\mathfrak{g}_\alpha
  \tag{6.7.54}
\end{equation}
设 $\Pi$ 为 $\mathfrak{g}$ 的单根之集。把根的高度函数记作 $\mathrm{ht}$，令 $\mathfrak{g}_i := \bigoplus_{\mathrm{ht}(\alpha)=i}\mathfrak{g}_\alpha$，则 $\mathfrak{g}=\bigoplus_{i\in\Z}\mathfrak{g}_i$ 是分次 Lie 代数。
\begin{equation}
  f = \sum_{\alpha\in\Pi}e_\alpha\in\mathfrak{g}_1\cong\mathfrak{g}^*_{-1},\qquad
  e = \sum_{\alpha\in\Pi}e_{-\alpha}\in\mathfrak{g}_{-1}\cong\mathfrak{g}^*_1
  \tag{6.7.55}
\end{equation}

% [原文 p.95]
此时，$e+f$ 的 $\ad^*_R$ 轨道之元可写成
\begin{equation}
  L = p + \sum_{\alpha\in\Pi}e^{(\alpha,q)}e_\alpha + e
  \tag{6.7.56}
\end{equation}
而 Hamilton 函数
\begin{equation}
  H(\xi) = \frac{1}{2}(\xi,\xi)
  \tag{6.7.57}
\end{equation}
表为
\begin{equation}
  H = \frac{1}{2}(p,p) + \sum_{\alpha\in\Pi}(e_\alpha,e_{-\alpha})e^{(\alpha,q)}
  \tag{6.7.58}
\end{equation}

本节中 Toda 格子的构造对应于 Lie 群的 Bruhat 分解。另一方面，6.3 节中的构造对应于 Lie 群的 Iwasawa 分解。
% [Tutor Note: 原书正文作「岩澤分解」；「岩澤」即 Iwasawa 的日文写法，故译作 Iwasawa 分解。]
这两种构造之间也可以建立联系，关于这一点例如可参看 [110, II.3.2]。

\chapterreview{%
本章从 Liouville--Arnold 定理出发，经 Lie 群的 Cartan、Iwasawa、Bruhat 三种分解，构造出完全可积的 Toda 格子，最后以 R 矩阵与 AKS 定理给出统一图景。
R 矩阵的本质是把 Lie 括号作 $\frac12(R\pm 1)$ 型变形：$R$ 为直和分解型 $P_+-P_-$ 时条件退化为 $[\mathfrak{g}_+,\mathfrak{g}_-]$ 的各向同性子空间结构，经典 Yang--Baxter 方程恰是 Jacobi 恒等式。
AKS 定理说明 Casimir 函数在 Lie--Poisson 结构下的 Hamilton 流必取 Lax 形式，从而第一积分自动对合，这是可积性的代数来源。
因子化定理进一步把 Lax 方程的解写成群上的因子化 $e^{tX_0}=g_+^{-1}g_-$，使流可用余伴随轨道上的 $g_+$ 表示。
最后对 $\mathfrak{sl}(n,\R)$ 取 $m=0,n=2$，轨道即为 $2(n-1)$ 维的 Toda 相空间，$a_i=e^{q_i-q_{i+1}}$、$b_i=p_i$ 给出 Darboux 坐标，Hamilton 函数正是经典 Toda 能量。
整套构造表明：完全可积性可由「R 矩阵 + 分次结构」这一纯代数数据编码，Bruhat 分解与 Iwasawa 分解分别给出 R 矩阵的两条实现路径。
}
